% EXEC_SUMMARY.tex — tw-finscope 執行摘要
% 編譯方式: xelatex EXEC_SUMMARY.tex
%
% 部署完成後，請將下方兩行的佔位符替換為實際網址，再重新編譯：
%   \newcommand{\demourl}{https://FILL\_IN\_YOUR\_VERCEL\_URL.vercel.app}
%   \newcommand{\githuburl}{https://github.com/yukinoshita045/tw-finscope}

\documentclass[11pt, a4paper]{article}

% ── 字型與 CJK ──────────────────────────────────────────────
\usepackage{fontspec}
\usepackage{xeCJK}
\setCJKmainfont{PingFang SC}[
  BoldFont   = PingFang SC Semibold,
  ItalicFont = PingFang SC Light
]
\setCJKsansfont{Heiti SC}
\setCJKmonofont{Heiti SC}
\setmainfont{Times New Roman}[BoldFont=Times New Roman Bold]
\setsansfont{Helvetica Neue}

% ── 版面 ────────────────────────────────────────────────────
\usepackage[a4paper, top=1.8cm, bottom=1.8cm, left=2.2cm, right=2.2cm]{geometry}
\usepackage{parskip}
\setlength{\parskip}{4pt}
\setlength{\parindent}{0pt}

% ── 套件 ────────────────────────────────────────────────────
\usepackage{booktabs}
\usepackage{tabularx}
\usepackage{array}
\usepackage{multirow}
\usepackage{xcolor}
\usepackage{hyperref}
\usepackage{enumitem}
\usepackage{titlesec}
\usepackage{fancyhdr}
\usepackage{graphicx}
\usepackage{tcolorbox}
\usepackage{microtype}
\usepackage{amsmath}

% ── 顏色定義 ─────────────────────────────────────────────────
\definecolor{primary}{HTML}{1a56db}
\definecolor{accent}{HTML}{0e9f6e}
\definecolor{lightgray}{HTML}{f3f4f6}
\definecolor{darkgray}{HTML}{374151}
\definecolor{urlcolor}{HTML}{1e40af}

% ── hyperref ─────────────────────────────────────────────────
\hypersetup{
  colorlinks = true,
  urlcolor   = urlcolor,
  linkcolor  = primary,
  pdfauthor  = {Tommy Lin},
  pdftitle   = {tw-finscope 執行摘要},
}

% ── 標題樣式 ──────────────────────────────────────────────────
\titleformat{\section}
  {\large\bfseries\color{primary}}
  {\thesection.}{0.5em}{}[\titlerule]
\titlespacing{\section}{0pt}{10pt}{4pt}

\titleformat{\subsection}
  {\normalsize\bfseries\color{darkgray}}
  {\thesubsection}{0.5em}{}
\titlespacing{\subsection}{0pt}{6pt}{2pt}

% ── 頁首頁尾 ──────────────────────────────────────────────────
\pagestyle{fancy}
\fancyhf{}
\renewcommand{\headrulewidth}{0.4pt}
\fancyhead[L]{\small\color{darkgray}\textbf{tw-finscope}\ \ 台灣上市公司財務健康儀表板}
\fancyhead[R]{\small\color{darkgray}執行摘要 · 2026}
\fancyfoot[C]{\small\thepage}

% ── 網址命令（部署後替換）────────────────────────────────────
\newcommand{\demourl}{https://FILL\_IN\_AFTER\_DEPLOY.vercel.app}
\newcommand{\githuburl}{https://github.com/yukinoshita045/tw-finscope}

% ══════════════════════════════════════════════════════════════
\begin{document}
% ══════════════════════════════════════════════════════════════

% ── 標題區塊 ──────────────────────────────────────────────────
\begin{tcolorbox}[
  colback=primary!8, colframe=primary, arc=3pt,
  boxrule=0.8pt, left=8pt, right=8pt, top=6pt, bottom=6pt
]
  {\LARGE\bfseries\color{primary} tw-finscope}\\[4pt]
  {\large\color{darkgray} 台灣上市公司財務健康儀表板}\\[6pt]
  \small
  \textbf{Demo：}
  \href{\demourl}{\textcolor{urlcolor}{\demourl}}
  \quad\textbf{GitHub：}
  \href{\githuburl}{\textcolor{urlcolor}{\githuburl}}
\end{tcolorbox}

\vspace{4pt}

% ── 一、問題與動機 ────────────────────────────────────────────
\section{問題與動機}

台灣公開資訊觀測站（MOPS）公開了所有上市公司財報，但資料分散在三種報表（資產負債
/損益/現金流量），欄位以中文命名、日期為民國紀年、金額為字串並含全形破折號與括號
負數，使一般使用者難以快速比較公司財務體質。\textbf{tw-finscope}
自動拉取並標準化這三份報表、計算十項衍生指標、寫入資料庫，以互動儀表板呈現，
讓使用者無需下載任何資料即可比較不同公司與產業的財務狀況。

% ── 二、資料工程亮點 ──────────────────────────────────────────
\section{資料工程亮點}

\textbf{E \;抽取：}Python \texttt{httpx} 呼叫 MOPS 三個 endpoint（\texttt{t164sb03/04/05}），
帶 browser-like headers 通過防爬機制；\texttt{tenacity} 三次指數回退重試，請求間 1.2 秒節流。

\textbf{T \;轉換：}
\begin{enumerate}[leftmargin=1.5em, itemsep=1pt, topsep=2pt]
  \item 民國年↔西元年雙向轉換
  \item 中文科目 → canonical key（\texttt{line\_items.py}，涵蓋 BS/IS/CF 共 60+ 對應）
  \item 字串清洗——千分位逗號、全形空白、括號表負、\texttt{－} 表 NA
  \item \textbf{計算 10 項衍生指標}——毛利率 / 營業利益率 / 淨利率、ROE、ROA、
        流動比率、負債比率、營業 FCF、營收 YoY、淨利 YoY
  \item 三個百分比指標自動偵測是否需要 $\times 100$ 標準化
\end{enumerate}

\textbf{L \;載入：}\texttt{ON CONFLICT DO UPDATE} 對 \texttt{(company\_id, year, season, statement\_type,
line\_item\_key)} 做冪等 upsert——任何時候重跑都不會產生重複列。
Schema 分三層：\texttt{statements}（長表原始）→ \texttt{metrics}（寬表黃金）→ \texttt{etl\_runs}（可追溯）。

\textbf{規模：}預設 40 家公司 × 16 季 = 640 個期間 × 3 份報表 $\approx$ 2,500 筆指標列。

% ── 三、視覺化 ───────────────────────────────────────────────
\section{視覺化}

四種互補的圖表類型，全程互動：

\begin{tabular}{@{}lp{12cm}@{}}
\toprule
\textbf{圖表} & \textbf{說明} \\
\midrule
損益趨勢線圖 &
  多公司營收/毛利/淨利時序對比，可疊加同產業平均（虛線）做基準\\[3pt]
利潤率分組長條圖 &
  毛利率/營業利益率/淨利率三層分解，看「賺錢的品質」\\[3pt]
財務比率雷達圖 &
  5 維（毛利率/營業利益率/淨利率/ROE/ROA）跨公司一眼比較體質\\[3pt]
產業排名表 &
  9 種指標自由切換，可按產業過濾，顯示前 25 名\\
\midrule
\multicolumn{2}{@{}l}{\textbf{財務健康分數卡（複合創新指標）}} \\
\multicolumn{2}{@{}p{14.5cm}}{
  把 5 項比率加權成 0–100 單一分數：
  毛利率（30\%）+ ROE（25\%）+ 流動比率（20\%）+ 低負債（15\%）+ 營收成長（10\%）；
  分數依區間以顏色（綠/藍/橘/紅）即時呈現，讓非財務背景使用者立即理解。
}\\
\bottomrule
\end{tabular}

% ── 四、資料更新機制 ──────────────────────────────────────────
\section{資料更新機制}

\texttt{.github/workflows/etl.yml}：cron 每週一台北時間 10:00 自動執行 ETL，寫入 Postgres，
同時支援 \texttt{workflow\_dispatch} 手動觸發（可選擇公司清單和年份）。
\texttt{etl\_runs} 表記錄每次起訖、處理公司數、寫入列數、狀態與註記，
構成完整 data lineage。儀表板右上角的「資料更新」徽章直接讀此表，
使用者隨時知道資料新鮮度。

% ── 五、系統架構 ──────────────────────────────────────────────
\section{系統架構}

\begin{center}
\small
\begin{tabular}{@{}l@{\ $\rightarrow$\ }l@{\ $\rightarrow$\ }l@{\ $\rightarrow$\ }l@{}}
GitHub Actions & Python ETL & PostgreSQL（Render） & FastAPI（Render）\\
（每週/手動） & extract→transform→load &  & /api/companies …\\
 &  & & $\downarrow$ JSON + CORS \\
 &  & & React + Chart.js（Vercel）\\
\end{tabular}
\end{center}

% ── 六、技術棧 ───────────────────────────────────────────────
\section{技術棧}

\begin{tabular}{@{}ll@{}}
\textbf{後端} & Python 3.12 · FastAPI · SQLAlchemy 2 · httpx · tenacity \\
\textbf{前端} & React 18 · Vite · Chart.js + react-chartjs-2 \\
\textbf{資料庫} & PostgreSQL 16（生產）/ SQLite（本地與離線 demo）\\
\textbf{DevOps} & GitHub Actions · Render（API + DB）· Vercel（前端）\\
\end{tabular}

% ── 七、連結 ─────────────────────────────────────────────────
\section{連結}

\begin{tabular}{@{}ll@{}}
\textbf{Demo 網址} &
  \href{\demourl}{\textcolor{urlcolor}{\demourl}} \\[4pt]
\textbf{Backend API Docs} &
  \href{https://tw-finscope-api.onrender.com/docs}{https://tw-finscope-api.onrender.com/docs} \\[4pt]
\textbf{GitHub Repository} &
  \href{\githuburl}{\textcolor{urlcolor}{\githuburl}} \\[4pt]
\textbf{GitHub Actions} &
  \href{\githuburl/actions}{\textcolor{urlcolor}{\githuburl/actions}} \\
\end{tabular}

\vspace{6pt}
\begin{tcolorbox}[
  colback=accent!8, colframe=accent, arc=3pt,
  boxrule=0.6pt, left=6pt, right=6pt, top=4pt, bottom=4pt
]
\small\color{darkgray}
\textbf{注意：}Demo URL（\texttt{FILL\_IN\_AFTER\_DEPLOY}）為部署前佔位符。
完成 Vercel 部署後，請替換第 42 行的 \texttt{\textbackslash demourl} 並重新編譯 PDF。
\end{tcolorbox}

\end{document}
